% !TeX program = xelatex
% Seguidor de línea con el Unitree H1-2: presentación de la socialización de los retos técnicos.
% Compilar con XeLaTeX, dos veces (las posiciones son absolutas): ./compilar.sh
\documentclass[aspectratio=169,10pt]{beamer}
\usepackage{fontspec}
\usepackage{xfp}
\usepackage{tikz}
\usetikzlibrary{arrows.meta,calc,positioning,shapes.geometric,decorations.pathreplacing}
\usepackage{graphicx}
\usepackage{array,colortbl}
\usepackage{ragged2e}
\usepackage{amsmath,amssymb}
% Página, fuentes y colores (base de las presentaciones de PI 2: 16:9, 960 × 540 bp).
% Coordenadas en bp (1 bp = 1/72 pulgada), con el origen arriba a la izquierda.
\geometry{paperwidth=960bp,paperheight=540bp}
\setbeamertemplate{navigation symbols}{}
\setbeamertemplate{headline}{}
\setbeamertemplate{footline}{}
\setbeamertemplate{background canvas}{}
\setbeamersize{text margin left=0pt,text margin right=0pt}
\setbeamercolor{normal text}{fg=black,bg=white}
\usefonttheme{professionalfonts}
\setlength{\parindent}{0pt}
\setlength{\parskip}{0pt}
\setlength{\emergencystretch}{1em}
\hyphenpenalty=10000
\exhyphenpenalty=10000

% Fuentes portátiles incluidas en fuentes/. Cambiar a \FuentesMicrosofttrue si Arial y
% Calibri están instaladas.
\newif\ifFuentesMicrosoft
\FuentesMicrosoftfalse
\ifFuentesMicrosoft
  \newfontfamily\FuenteArial{Arial}
  \newfontfamily\FuenteCalibri{Calibri}
  \setsansfont{Calibri}
\else
  \newfontfamily\FuenteArial{LiberationSans}[Path=fuentes/,Extension=.ttf,
    UprightFont=*-Regular,BoldFont=*-Bold,ItalicFont=*-Italic,BoldItalicFont=*-BoldItalic]
  \newfontfamily\FuenteCalibri{Carlito}[Path=fuentes/,Extension=.ttf,
    UprightFont=*-Regular,BoldFont=*-Bold,ItalicFont=*-Italic,BoldItalicFont=*-BoldItalic]
  \setsansfont{Carlito}[Path=fuentes/,Extension=.ttf,
    UprightFont=*-Regular,BoldFont=*-Bold,ItalicFont=*-Italic,BoldItalicFont=*-BoldItalic]
\fi
\newfontfamily\FuenteSpace{SpaceGrotesk}[Path=fuentes/,Extension=.ttf,
  UprightFont=*-Regular,BoldFont=*-Bold]

% Colores UTEC (los de la presentación de 5 minutos).
\definecolor{GrisUTEC}{HTML}{4E4B4C}
\definecolor{CianUTEC}{HTML}{00E1FF}
\definecolor{GrisClaroUTEC}{HTML}{E4E4E4}
% Derivados para filas alternas, reglas y trazos finos sobre blanco.
\definecolor{GrisFila}{HTML}{F2F2F2}
\definecolor{GrisRegla}{HTML}{C8C8C8}
\definecolor{GrisMedio}{HTML}{8A8788}

% Comandos de la presentación del seguidor de línea, sobre estilo.tex.
% Formato de la presentación de resultados de PI 2 (cajas con posición fija, tablas con
% cabecera y filas alternas, pies «Fig.»), con los colores y el logo de UTEC.
% Todo va dentro de un Lienzo: coordenadas en bp, origen arriba a la izquierda, y hacia abajo.
%
% Zonas fijas de las diapositivas de contenido (las comprueba verificar_margenes.py):
%   logo x 36–96, y 14–52 · título desde x 110 · etiqueta del bloque arriba a la derecha
%   regla y 64–68 · contenido y 80–505 · pie de página y 515–530

\newenvironment{Lienzo}{%
  \begin{tikzpicture}[remember picture,overlay]
  \begin{scope}[shift={(current page.north west)},x=1bp,y=-1bp]
}{%
  \end{scope}\end{tikzpicture}%
}

% ── Registro para verificar_margenes.py ──────────────────────────────────
\newwrite\RegistroLayout
\immediate\openout\RegistroLayout=\jobname.layout
% Imagen: x, y (borde superior), ancho, alto, ruta.
\newcommand{\Imagen}[5]{%
  \immediate\write\RegistroLayout{\thepage;#1;#2;#3;#4;#5}%
  \node[anchor=north west,inner sep=0pt,outer sep=0pt] at (#1,#2)
    {\includegraphics[width=#3bp,height=#4bp]{#5}};}
\newcommand{\Figura}[5]{\Imagen{#1}{#2}{#3}{#4}{graficas/#5.pdf}}
\newcommand{\Foto}[5]{\Imagen{#1}{#2}{#3}{#4}{imagenes/fotos/#5.jpg}}
\newcommand{\Doc}[5]{\Imagen{#1}{#2}{#3}{#4}{imagenes/docs/#5}}

% ── Texto ────────────────────────────────────────────────────────────────
% Texto de una línea: x, línea base, fuente (Arial, Calibri, Space), tamaño, color HTML, contenido.
\newcommand{\Texto}[6]{%
  \node[anchor=base west,inner sep=0pt,outer sep=0pt,
    font=\csname Fuente#3\endcsname\fontsize{#4bp}{#4bp}\selectfont]
    at (#1,#2) {\color[HTML]{#5}#6};}
% Párrafo: x, línea base de la 1.ª línea, ancho, fuente, tamaño, interlineado, color, contenido.
\newcommand{\Parrafo}[8]{%
  \node[anchor=base west,inner sep=0pt,outer sep=0pt] at (#1,#2)
    {\begin{minipage}[t]{#3bp}
      \csname Fuente#4\endcsname\fontsize{#5bp}{#6bp}\selectfont\color[HTML]{#7}%
      \RaggedRight\hyphenpenalty=10000\exhyphenpenalty=10000\setlength{\parindent}{0pt}#8
    \end{minipage}};}
\newcommand{\ParrafoC}[8]{%
  \node[anchor=base west,inner sep=0pt,outer sep=0pt] at (#1,#2)
    {\begin{minipage}[t]{#3bp}
      \csname Fuente#4\endcsname\fontsize{#5bp}{#6bp}\selectfont\color[HTML]{#7}%
      \centering\hyphenpenalty=10000\exhyphenpenalty=10000\setlength{\parindent}{0pt}#8
    \end{minipage}};}
% Cuerpo en Arial negro (como el texto de la presentación de resultados): x, y, ancho, tamaño, contenido.
% Dentro: \Sub{subtítulo}, \Vin{viñeta}, \Esp{bp}.
\newcommand{\Cuerpo}[5]{%
  \Parrafo{#1}{#2}{#3}{Arial}{#4}{\fpeval{1.22*#4}}{000000}{#5}}
\newcommand{\CuerpoC}[5]{%
  \Parrafo{#1}{#2}{#3}{Calibri}{#4}{\fpeval{1.2*#4}}{000000}{#5}}
\newcommand{\Vin}[1]{\par{\leftskip=17bp\relax\parindent=-17bp\relax
  \makebox[17bp][l]{\color{CianUTEC}\rule[0.18em]{0.48em}{0.48em}}#1\par}}
\newcommand{\VinG}[1]{\par{\leftskip=17bp\relax\parindent=-17bp\relax
  \makebox[17bp][l]{\color{GrisUTEC}\rule[0.18em]{0.48em}{0.48em}}#1\par}}
\newcommand{\Sub}[1]{\par{\bfseries #1\par}}
\newcommand{\Esp}[1]{\par\vspace{#1bp}}
% Resaltado en el texto: negrita negra sobre una franja cian.
\newcommand{\Marca}[1]{{\setlength{\fboxsep}{1.2bp}\colorbox{CianUTEC}{\strut\bfseries #1}}}

% Pie de figura y de tabla: «Fig.»/«Tabla.» en negrita, Calibri gris. x, y, ancho, texto.
\newcommand{\Pie}[4]{\Parrafo{#1}{#2}{#3}{Calibri}{13}{15.5}{4E4B4C}{\textbf{Fig.}~#4}}
\newcommand{\PieTabla}[4]{\Parrafo{#1}{#2}{#3}{Calibri}{13}{15.5}{4E4B4C}{\textbf{Tabla.}~#4}}
\newcommand{\Nota}[4]{\Parrafo{#1}{#2}{#3}{Calibri}{12}{14.5}{4E4B4C}{#4}}

% ── Cajas y cifras ────────────────────────────────────────────────────────
% Recuadro gris claro: x, y, ancho, alto.
\newcommand{\Recuadro}[4]{\fill[GrisClaroUTEC!55,rounded corners=3bp] (#1,#2) rectangle ({#1+#3},{#2+#4});}
% Barra cian de acento a la izquierda de un bloque: x, y superior, y inferior.
\newcommand{\Acento}[3]{\fill[CianUTEC] (#1,#2) rectangle ({#1+5},#3);}
% Tarjeta: recuadro gris claro con banda cian arriba y título. x, y, ancho, alto, título.
\newcommand{\Tarjeta}[5]{%
  \fill[GrisClaroUTEC!55,rounded corners=3bp] (#1,#2) rectangle ({#1+#3},{#2+#4});
  \fill[CianUTEC] (#1,#2) rectangle ({#1+#3},{#2+5});
  \Texto{#1+12}{#2+27}{Calibri}{18}{000000}{\bfseries #5}}
% Cabecera oscura con texto blanco: x, y, ancho, alto, texto.
\newcommand{\Banda}[5]{%
  \fill[GrisUTEC] (#1,#2) rectangle ({#1+#3},{#2+#4});
  \node[anchor=west,inner sep=0pt,outer sep=0pt,text=white,
    font=\FuenteCalibri\bfseries\fontsize{16bp}{16bp}\selectfont] at ({#1+10},{#2+#4/2}) {#5};}
% Cifra con su explicación debajo y la barra cian a la izquierda:
% [tamaño]{x}{y borde superior}{ancho}{cifra}{explicación}
\newcommand{\CifraBarra}[6][32]{%
  \fill[CianUTEC] ({#2-13},{#3+2}) rectangle ++(5,{#1*0.95});
  \node[anchor=north west,inner sep=0pt,outer sep=0pt] at (#2,#3)
    {\begin{minipage}[t]{#4bp}\setlength{\parindent}{0pt}\RaggedRight
      \hyphenpenalty=10000\exhyphenpenalty=10000\relax
      {\FuenteCalibri\bfseries\fontsize{#1bp}{\dimexpr#1bp*11/10\relax}\selectfont
        \color{black}#5\par}\vspace{2bp}
      {\FuenteCalibri\fontsize{15bp}{18bp}\selectfont\color{GrisUTEC}#6\par}
    \end{minipage}};}
% Número en un cuadro cian (pasos, etapas): x, y (borde superior), lado, número.
\newcommand{\Numero}[4]{%
  \node[anchor=north west,fill=CianUTEC,minimum width=#3bp,minimum height=#3bp,
    inner sep=0pt,outer sep=0pt,text=black,
    font=\FuenteCalibri\bfseries\fontsize{\fpeval{0.6*#3}bp}{\fpeval{0.6*#3}bp}\selectfont] at (#1,#2) {#4};}
% Etiqueta (chip) cian: x, y (borde superior), texto.
\newcommand{\Chip}[3]{%
  \node[anchor=north west,fill=CianUTEC,inner xsep=6bp,inner ysep=4bp,outer sep=0pt,text=black,
    font=\FuenteCalibri\bfseries\fontsize{12bp}{12bp}\selectfont] at (#1,#2) {#3};}
\newcommand{\ChipG}[3]{%
  \node[anchor=north west,fill=GrisUTEC,inner xsep=6bp,inner ysep=4bp,outer sep=0pt,text=white,
    font=\FuenteCalibri\bfseries\fontsize{12bp}{12bp}\selectfont] at (#1,#2) {#3};}

% Botón que abre un vídeo con el reproductor del sistema: x, y (borde superior), archivo, texto.
% El enlace es una anotación /Launch del tamaño del botón (xdvipdfmx solo hace clicable el texto de un
% \href), así que todo el botón abre el vídeo.
\newcommand{\BotonVideo}[4]{%
  \fill[CianUTEC,rounded corners=4bp] (#1,#2) rectangle ({#1+150},{#2+26});
  \fill[black] ({#1+11},{#2+6}) -- ({#1+11},{#2+20}) -- ({#1+23},{#2+13}) -- cycle;
  \node[anchor=west,inner sep=0pt,outer sep=0pt,text=black,
    font=\FuenteCalibri\bfseries\fontsize{14bp}{14bp}\selectfont] at ({#1+31},{#2+13}) {#4};
  \node[anchor=north west,inner sep=0pt,outer sep=0pt] at (#1,#2)
    {\hbox to 150bp{\rule{0pt}{26bp}\special{pdf:ann width 150bp height 26bp depth 0bp
      << /Type /Annot /Subtype /Link /Border [0 0 0] /A << /S /Launch /F (videos/#3) >> >>}\hfil}};}

% ── Tablas: cabecera gris UTEC con texto blanco, filas alternas y reglas grises ──
% \Tabla{x}{y}{columnas}{filas}; filas con \FCab/\FGris/\FBlanca y \Cab/\Cel/\CelI/\CelN/\CelNI
\newcommand{\BloqueSuperior}[3]{\node[anchor=north west,inner sep=0pt,outer sep=0pt] at (#1,#2) {#3};}
\newcommand{\Celda}[6]{%
  \parbox[c][#2bp][c]{#1bp}{%
    \hspace*{6bp}\begin{minipage}{\dimexpr#1bp-12bp\relax}%
      \FuenteCalibri\fontsize{#3bp}{\dimexpr#3bp*6/5\relax}\selectfont#5%
      \ifx c#4\centering\else\raggedright\fi
      #6\par
    \end{minipage}}}
\newcommand{\CeldaR}[6]{\Celda{#1}{#2}{#3}{#4}{#5}{\strut#6\strut}}
\newcommand{\Tabla}[4]{%
  \BloqueSuperior{#1}{#2}{%
    \begingroup\FuenteCalibri\setlength{\tabcolsep}{0pt}%
    \setlength{\arrayrulewidth}{0.5bp}\arrayrulecolor{GrisRegla}%
    \begin{tabular}{@{}#3@{}}#4\end{tabular}\endgroup}}
\newcommand{\FCab}{\rowcolor{GrisUTEC}}
\newcommand{\FGris}{\rowcolor{GrisFila}}
\newcommand{\FBlanca}{\rowcolor{white}}
\newcommand{\FCian}{\rowcolor{CianUTEC!35}}
\newcommand{\Cab}[3]{\CeldaR{#1}{#2}{#3}{c}{\bfseries\color{white}}}
\newcommand{\CabI}[3]{\CeldaR{#1}{#2}{#3}{l}{\bfseries\color{white}}}
\newcommand{\Cel}[3]{\CeldaR{#1}{#2}{#3}{c}{\mdseries\color{black}}}
\newcommand{\CelI}[3]{\CeldaR{#1}{#2}{#3}{l}{\mdseries\color{black}}}
\newcommand{\CelN}[3]{\CeldaR{#1}{#2}{#3}{c}{\bfseries\color{black}}}
\newcommand{\CelNI}[3]{\CeldaR{#1}{#2}{#3}{l}{\bfseries\color{black}}}

% ── Diagramas ────────────────────────────────────────────────────────────
\tikzset{
  bloque/.style={draw=GrisUTEC,line width=1bp,fill=white,rounded corners=3bp,align=center,
    inner sep=4bp,font=\FuenteCalibri\fontsize{14bp}{16bp}\selectfont,text=black},
  bloqueC/.style={bloque,fill=CianUTEC,draw=CianUTEC},
  bloqueG/.style={bloque,fill=GrisUTEC,draw=GrisUTEC,text=white},
  bloqueL/.style={bloque,fill=GrisClaroUTEC!55,draw=GrisClaroUTEC},
  flecha/.style={-{Stealth[length=7bp,width=6bp]},line width=1.3bp,color=GrisUTEC},
  flechaC/.style={-{Stealth[length=7bp,width=6bp]},line width=1.6bp,color=black},
  rotulo/.style={font=\FuenteCalibri\fontsize{12bp}{14bp}\selectfont,text=GrisUTEC,align=center,inner sep=2bp},
}

% ── Diapositiva de contenido: cabecera, regla y pie ──────────────────────
\newcommand{\PieDePagina}{Socialización de los retos técnicos – Unitree H1-2 · Seguidor de línea · 13 de noviembre de 2026}
\newcommand{\NumPag}{%
  \node[anchor=base east,inner sep=0pt,outer sep=0pt,
    font=\FuenteCalibri\fontsize{13bp}{13bp}\selectfont] at (924,526)
    {\color{GrisUTEC}\insertframenumber};}
% \Cabecera{bloque}{título}: logo, título, etiqueta del bloque, regla y pie.
\newcommand{\Cabecera}[2]{%
  \Imagen{36}{14}{60}{37.9}{imagenes/utec_logo1.png}
  \node[anchor=base west,inner sep=0pt,outer sep=0pt,
    font=\FuenteCalibri\bfseries\fontsize{30bp}{30bp}\selectfont] at (110,47) {\color{black}#2};
  \node[anchor=north east,fill=CianUTEC,inner xsep=6bp,inner ysep=4bp,outer sep=0pt,text=black,
    font=\FuenteCalibri\bfseries\fontsize{11bp}{11bp}\selectfont] at (924,20) {\MakeUppercase{#1}};
  \fill[GrisRegla] (36,65.6) rectangle (924,66.4);
  \fill[CianUTEC] (36,63.5) rectangle (96,68.5);
  \node[anchor=base west,inner sep=0pt,outer sep=0pt,
    font=\FuenteCalibri\fontsize{11bp}{11bp}\selectfont] at (36,526) {\color{GrisUTEC}\PieDePagina};
  \NumPag}

% Fila «problema → cambio → resultado» de las diapositivas de resultados:
% x, y (borde superior), ancho, alto, tipo (PROBLEMA, CAMBIO o RESULTADO), texto.
\newcommand{\Fila}[6]{%
  \fill[GrisClaroUTEC!55,rounded corners=3bp] (#1,#2) rectangle ({#1+#3},{#2+#4});
  \ifnum\strcmp{#5}{RESULTADO}=0
    \fill[CianUTEC,rounded corners=3bp] (#1,#2) rectangle ({#1+108},{#2+#4});
    \node[anchor=center,inner sep=0pt,text=black,font=\FuenteCalibri\bfseries\fontsize{13bp}{13bp}\selectfont]
      at ({#1+54},{#2+#4/2}) {#5};
  \else
    \fill[GrisUTEC,rounded corners=3bp] (#1,#2) rectangle ({#1+108},{#2+#4});
    \node[anchor=center,inner sep=0pt,text=white,font=\FuenteCalibri\bfseries\fontsize{13bp}{13bp}\selectfont]
      at ({#1+54},{#2+#4/2}) {#5};
  \fi
  \node[anchor=west,inner sep=0pt,outer sep=0pt] at ({#1+120},{#2+#4/2})
    {\begin{minipage}{\dimexpr#3bp-130bp\relax}\FuenteCalibri\fontsize{14.5bp}{17.5bp}\selectfont
      \RaggedRight\hyphenpenalty=10000\exhyphenpenalty=10000 #6\end{minipage}};}

\hypersetup{pdftitle={Seguimiento autónomo de línea con el humanoide Unitree H1-2},
  pdfauthor={Miguel Olortegui, Sergio Morales, Axel Cespedes},
  pdfsubject={Socialización de los retos técnicos - Unitree H1-2}}
\begin{document}
% ── Apertura
% Portada: logo, evento, título, equipo y fecha; a la derecha, el H1-2 en la recta del nivel 1.
\begin{frame}[plain]
\begin{Lienzo}
\Foto{568}{0}{392}{540}{portada}
\fill[CianUTEC] (560,0) rectangle (568,540);
\Imagen{48}{34}{114}{72}{imagenes/utec_logo1.png}
\Chip{48}{136}{SOCIALIZACIÓN DE LOS RETOS TÉCNICOS – UNITREE H1-2}
\Parrafo{48}{200}{490}{Calibri}{37}{41}{000000}{\bfseries Seguimiento autónomo de línea con el humanoide Unitree H1-2}
\Parrafo{48}{330}{470}{Calibri}{19}{23}{4E4B4C}{Reto técnico: el H1-2 como seguidor de línea con realimentación por la cámara RealSense D435i de la cabeza}
\Acento{48}{388}{470}
\Parrafo{64}{402}{460}{Calibri}{17}{22.5}{000000}{\textbf{Miguel Olortegui} · Operador de seguridad\par
\textbf{Sergio Morales} · Responsable de desarrollo\par
\textbf{Axel Cespedes} · Responsable de documentación}
\Texto{48}{506}{Calibri}{16}{4E4B4C}{13 de noviembre de 2026}
\end{Lienzo}
\end{frame}

% Equipo y responsabilidades: los tres roles del reto, con las tareas de cada uno.
\begin{frame}[plain]
\begin{Lienzo}
\Cabecera{Equipo}{Equipo y responsabilidades}
% Tres tarjetas: rol (banda gris), nombre y tareas.
\foreach \x/\rol/\nombre in {36/Operador de seguridad/Miguel Olortegui,
                            337/Responsable de desarrollo/Sergio Morales,
                            638/Responsable de documentación/Axel Cespedes}{
  \fill[GrisClaroUTEC!55,rounded corners=3bp] (\x,84) rectangle ({\x+286},444);
  \fill[GrisUTEC] (\x,84) rectangle ({\x+286},116);
  \node[anchor=west,inner sep=0pt,text=white,font=\FuenteCalibri\bfseries\fontsize{16bp}{16bp}\selectfont] at ({\x+12},100) {\rol};
  \fill[CianUTEC] (\x,116) rectangle ({\x+286},120);
  \Texto{\x+12}{146}{Calibri}{21}{000000}{\bfseries\nombre}
}
\Cuerpo{48}{176}{264}{13.5}{%
\Vin{Mando en la mano en todas las tiradas, con L2\,+\,B (parada de emergencia) listo y sin otra tarea.}\Esp{5}
\Vin{Antes de cada tirada: FSM 201 (L2\,+\,arriba), motores sin fallo, batería cargada y el robot en el cuadro de salida.}\Esp{5}
\Vin{Franja de 1.5 m libre alrededor de la pista; detiene la tirada ante cualquier riesgo.}\Esp{5}
\Vin{Lleva el robot en el simulacro y en la grabación de los conjuntos de datos.}}
\Cuerpo{349}{176}{264}{13.5}{%
\Vin{Arquitectura de cinco bloques y su código: percepción, estimación, control, supervisor y registro.}\Esp{5}
\Vin{Calibración de la cámara, caracterización de la marcha, simulador y pruebas automáticas.}\Esp{5}
\Vin{Despliegue en el PC2, simulacros y tiradas; analiza cada registro antes de repetir.}\Esp{5}
\Vin{Ajustes por nivel: giro inicial, ajuste local, pausas, búsqueda en arco, parada adelantada.}}
\Cuerpo{650}{176}{264}{13.5}{%
\Vin{Registro de cada tirada: nivel, escala, nota, resultado e intervenciones.}\Esp{5}
\Vin{Medidas con cinta en la pista: dónde para respecto de la barra y dónde quedan los pies tras la esquina.}\Esp{5}
\Vin{Vídeos y fotos de los cuatro niveles; conjuntos de datos y evidencias ordenados.}\Esp{5}
\Vin{Métricas de la sección 9 e informe con las respuestas 6.1–6.4.}}
% Reglas comunes
\Recuadro{36}{454}{888}{52}
\Acento{36}{454}{506}
\Parrafo{52}{474}{860}{Calibri}{15}{18}{000000}{\textbf{Reglas comunes (secciones 8 y 10):} ninguna orden autónoma antes del Hito 3; las primeras tiradas de cada nivel, a media escala; una tirada fallida se analiza con su registro antes de repetirla.}
\end{Lienzo}
\end{frame}

% ── I · Interpretación del reto
% Interpretación del reto (secciones 1 y 2 del enunciado) y por qué no es un seguidor de ruedas.
\begin{frame}[plain]
\begin{Lienzo}
\Cabecera{I · Interpretación del reto}{Qué pide el reto y por qué no es trivial}
\Cuerpo{36}{100}{410}{15}{%
\Sub{Qué pide (sección 1)}\Esp{3}
Que el H1-2 camine siguiendo una línea del suelo, con la RealSense D435i de la cabeza como \textbf{única realimentación de posición}, y que se detenga solo en la barra de fin.\Esp{9}
\Sub{Qué escribe el equipo}\Esp{3}
No el control de las piernas: la marcha es la de Unitree. El equipo escribe el lazo que decide qué pedir a \texttt{LocoClient.Move(vx, vy, vyaw)} a partir de la imagen y de la IMU.\Esp{9}
\Sub{Objetivo general (sección 2)}\Esp{3}
Seguimiento en lazo cerrado con percepción en la D435i, control sobre LocoClient y supervisión de seguridad, hasta el nivel más alto posible y con registro medible de cada tirada.}
\Recuadro{36}{410}{420}{94}
\Acento{36}{410}{504}
\Parrafo{52}{432}{396}{Calibri}{16}{19.5}{000000}{\textbf{Nuestra lectura:} medir primero el robot real (cámara y marcha) y diseñar después; la percepción, sobre datos grabados; nada anda solo antes del simulacro.}
% Columna derecha: dificultades, con lo medido.
\Texto{486}{100}{Arial}{15}{000000}{\bfseries Por qué no es un seguidor de ruedas (y lo que medimos)}
\foreach \i/\y/\tit/\txt in {
  1/118/{Cámara en la cabeza, sin cuello}/{A 1.66 m del suelo e inclinada 54°: hay que saber qué suelo se ve. Medido: de 0.10 a 4.07 m por delante.},
  2/195/{La imagen se balancea}/{Cadencia de 1.33–1.46 Hz: la línea oscila \textasciitilde1 cm y \textasciitilde1° en cada paso aunque el robot vaya recto.},
  3/272/{Retardo y deriva}/{La marcha responde con 0.3–0.45 s de retardo efectivo; el rumbo deriva \textasciitilde1.5°/s a la izquierda.},
  4/349/{Sin odometría}/{El odómetro del PC2 devuelve (0, 0, 0): la distancia solo sale de la cámara o de estimarla.},
  5/426/{Un humanoide de 70 kg}/{Un error de control no es salirse de la pista: es una caída. La seguridad va primero.}}{
  \Numero{486}{\y}{30}{\i}
  \Parrafo{528}{\y+13}{396}{Calibri}{17}{20}{000000}{\textbf{\tit}}
  \Parrafo{528}{\y+34}{396}{Calibri}{15}{18}{4E4B4C}{\txt}
}
\end{Lienzo}
\end{frame}

% Niveles (sección 4) y criterio de éxito (sección 9): la pista del PDF y la que se montó.
\begin{frame}[plain]
\begin{Lienzo}
\Cabecera{I · Interpretación del reto}{Cuatro niveles, cuatro pistas}
\tikzset{cinta/.style={line width=2.6bp,color=GrisUTEC},
         salida/.style={draw=GrisMedio,line width=0.6bp,dashed},
         barra/.style={line width=2.6bp,color=black}}
% Columnas: chip, qué pone a prueba, dibujo (vista desde arriba, 35 bp/m) y pista del PDF frente a la montada.
\foreach \x/\n/\prueba in {36/1/{Corregir un error inicial de orientación y la deriva; parar sin pasarse más de 30 cm.},
  262/2/{Seguir la curvatura con la zona ciega: la curva se vio varios pasos antes.},
  488/3/{Pérdida momentánea de la línea, cambio de luz y curvatura que cambia de signo.},
  714/4/{Una maniobra discontinua: detectar la esquina y girar sobre el eje.}}{
  \Chip{\x}{82}{NIVEL \n}
  \Parrafo{\x}{122}{210}{Calibri}{14}{16.5}{4E4B4C}{\prueba}
}
\Chip{776}{82}{MÉRITO}
% Nivel 1: recta de 4 m, salida a 10°
\begin{scope}[shift={(130,328)},x=35bp,y=35bp]
  \draw[salida] (-0.3,-0.45) rectangle (0.3,0.15);
  \draw[cinta] (0,0) -- (0,4);
  \draw[barra] (-0.3,4) -- (0.3,4);
  \draw[-{Stealth[length=6bp]},line width=1.6bp,color=CianUTEC] (0,-0.15) -- ++(100:0.9);
  \node[rotulo,anchor=west] at (0.15,0.35) {10°};
\end{scope}
% Nivel 2: recta 2 m, curva de 90° de radio 1 m a la derecha, recta 2 m
\begin{scope}[shift={(302,328)},x=35bp,y=35bp]
  \draw[salida] (-0.3,-0.45) rectangle (0.3,0.15);
  \draw[cinta] (0,0) -- (0,2) arc[start angle=180,end angle=90,radius=1] -- (3,3);
  \draw[barra] (3,2.7) -- (3,3.3);
  \node[rotulo,anchor=north] at (1.25,2.05) {R 1 m};
\end{scope}
% Nivel 3 (equipo): recta 1 m, semicírculo a la derecha y semicírculo a la izquierda de R 1 m, recta 1 m
\begin{scope}[shift={(508,328)},x=35bp,y=35bp]
  \fill[GrisUTEC,opacity=0.18] (3.45,1.05) rectangle (4.55,2.35);
  \draw[salida] (-0.3,-0.45) rectangle (0.3,0.15);
  \draw[cinta] (0,0) -- (0,1) arc[start angle=180,end angle=0,radius=1]
                arc[start angle=180,end angle=360,radius=1] -- (4,2);
  \draw[barra] (3.7,2) -- (4.3,2);
  \fill[white] (1.88,0.79) rectangle (2.12,1.21);
  \draw[GrisUTEC,line width=0.6bp] (1.9,0.95) -- (1.55,0.5);
  \node[rotulo,anchor=north] at (1.08,0.45) {hueco 40 cm};
  \node[rotulo,anchor=east] at (3.42,1.75) {sombra};
  \node[rotulo,anchor=north] at (1,1.9) {R 1 m};
\end{scope}
% Nivel 4: recta 2 m, esquina de 90° a la derecha, recta 2 m
\begin{scope}[shift={(760,328)},x=35bp,y=35bp]
  \draw[salida] (-0.3,-0.45) rectangle (0.3,0.15);
  \draw[cinta] (0,0) -- (0,2) -- (2,2);
  \draw[barra] (2,1.7) -- (2,2.3);
  \node[rotulo,anchor=north west] at (0.08,1.95) {90°};
  % escala
  \draw[line width=1.4bp] (1.1,0.05) -- (2.1,0.05);
  \node[rotulo,anchor=north] at (1.6,0.0) {1 m};
\end{scope}
% Pista del PDF frente a la montada
\foreach \x/\pdf/\nuestra in {36/{Recta de 4 m, salida desalineada 10° y barra.}/{Igual; salidas casi alineadas y a ±10°.},
  262/{Recta 2 m, curva de 90° de radio 2 m, recta 2 m.}/{Curva de radio \textasciitilde1 m a la derecha (por espacio).},
  488/{S con radios de 1.2 m, 40 cm sin línea y 1 m en sombra.}/{Dos semicírculos de 1 m: giran 180° a cada lado. Hueco en la unión; sombra apagando luces.},
  714/{Esquina en ángulo recto, 2 m de recta y barra.}/{Igual, con la esquina a la derecha.}}{
  \Parrafo{\x}{356}{210}{Calibri}{14}{16.5}{000000}{\textbf{PDF:} \pdf}
  \Parrafo{\x}{396}{210}{Calibri}{14}{16.5}{000000}{\textbf{Montada:} \nuestra}
}
% Criterio de éxito
\Recuadro{36}{448}{888}{58}
\Acento{36}{448}{506}
\Parrafo{52}{469}{864}{Calibri}{15}{18}{000000}{\textbf{Éxito (sección 9):} el robot llega a la barra sin intervención y sin que el pie salga más de 30 cm de la línea.}
\end{Lienzo}
\end{frame}

% ── II · Arquitectura propuesta (sección 5)
% Arquitectura mínima exigida (sección 5): los cinco bloques, sus interfaces y cómo se cumple cada requisito.
\begin{frame}[plain]
\begin{Lienzo}
\Cabecera{II · Arquitectura (sección 5)}{Cinco bloques con interfaces declaradas}
% Diagrama de flujo
\node[bloqueL,minimum width=96bp,minimum height=46bp] (cam) at (86,138) {D435i\\[-1bp]{\fontsize{11bp}{12bp}\selectfont IR 640×480, 30 fps}};
\node[bloqueC,minimum width=112bp,minimum height=46bp] (per) at (238,138) {\bfseries Percepción};
\node[bloqueC,minimum width=112bp,minimum height=46bp] (est) at (414,138) {\bfseries Estimación};
\node[bloqueC,minimum width=112bp,minimum height=46bp] (con) at (590,138) {\bfseries Control};
\node[bloqueG,minimum width=118bp,minimum height=46bp] (sup) at (770,138) {\bfseries Supervisor\\[-1bp]{\fontsize{11bp}{12bp}\selectfont + vigilante}};
\node[bloqueL,minimum width=150bp,minimum height=34bp] (imu) at (414,222) {IMU y salud del robot\\[-1bp]{\fontsize{11bp}{12bp}\selectfont \texttt{rt/lowstate}, 100 Hz}};
\node[bloqueL,minimum width=150bp,minimum height=34bp] (mar) at (770,222) {Marcha de Unitree\\[-1bp]{\fontsize{11bp}{12bp}\selectfont \texttt{Move} / \texttt{StopMove}}};
\node[bloque,dashed,minimum width=250bp,minimum height=30bp] (reg) at (172,222) {\bfseries Registro\ {\mdseries\fontsize{12bp}{12bp}\selectfont → CSV + resumen JSON}};
\draw[flecha] (cam) -- (per);
\draw[flecha] (per) -- node[rotulo,above=1bp] {medida} (est);
\draw[flecha] (est) -- node[rotulo,above=1bp] {estado} (con);
\draw[flecha] (con) -- node[rotulo,above=1bp] {orden} (sup);
\draw[flecha] (imu) -- (est);
\draw[flecha] (sup) -- node[rotulo,right=2bp] {20 Hz} (mar);
\draw[flecha] ($(imu.north east)+(-10,0)$) -- ++(0,-20) -| node[rotulo,pos=0.25,above=0bp] {salud} ($(sup.south west)+(16,0)$);
\draw[flecha,dashed] (per.south) -- (per.south |- reg.north);
\Nota{36}{94}{300}{Fotograma → medida → estado → orden → \texttt{Move}}
% Tabla de interfaces
\Tabla{36}{256}{p{104bp}p{266bp}p{242bp}p{276bp}}{%
\FCab
\CabI{104}{26}{13.5}{Bloque} & \CabI{266}{26}{13.5}{Entrada → salida} & \CabI{242}{26}{13.5}{Requisito del enunciado} & \CabI{276}{26}{13.5}{Cómo se cumple}\\\hline
\FGris
\CelNI{104}{43}{13}{Percepción} & \CelI{266}{43}{13}{Fotograma → desplazamiento (m), ángulo (rad), curvatura, confianza 0–1, tiempo} & \CelI{242}{43}{13}{Igual sobre el conjunto de datos que en vivo, sin cambiar el código} & \CelI{276}{43}{13}{La fuente (ZMQ o dataset) se elige fuera; 13 pruebas sintéticas y una de regresión}\\\hline
\FBlanca
\CelNI{104}{43}{13}{Estimación} & \CelI{266}{43}{13}{Medidas + IMU → estado de la línea respecto del robot} & \CelI{242}{43}{13}{Declarar el trato del balanceo y qué hace si cae la confianza} & \CelI{276}{43}{13}{Sin corregir roll/pitch por paso; yaw de la IMU. Sin medida: estima de avance}\\\hline
\FGris
\CelNI{104}{43}{13}{Control} & \CelI{266}{43}{13}{Estado → vx, vy, vyaw saturadas} & \CelI{242}{43}{13}{Ganancias y límites en un fichero, no en el código} & \CelI{276}{43}{13}{\texttt{config/seguidor.yaml}; topes 0.4 m/s, 0.2 m/s y 0.5 rad/s}\\\hline
\FBlanca
\CelNI{104}{43}{13}{Supervisor} & \CelI{266}{43}{13}{Estado + orden + salud del robot → orden final o parada} & \CelI{242}{43}{13}{Máquina de estados explícita; único que llama a Move y StopMove} & \CelI{276}{43}{13}{Seis estados; el vigilante mira cámara, IMU, FSM, motores e inclinación}\\\hline
\FGris
\CelNI{104}{43}{13}{Registro} & \CelI{266}{43}{13}{Todo lo anterior → CSV y resumen JSON por tirada} & \CelI{242}{43}{13}{Bastar para reconstruir cada tirada después} & \CelI{276}{43}{13}{Fotogramas, medidas, órdenes, \texttt{lowstate}, imágenes IR y \texttt{resumen.json}}\\\hline
}
\end{Lienzo}
\end{frame}

% Despliegue: procesos en el PC2 del robot y por qué el lazo no pasa por el portátil.
\begin{frame}[plain]
\begin{Lienzo}
\Cabecera{II · Arquitectura (sección 5)}{Despliegue: todo el lazo corre en el robot}
% Caja del PC2
\fill[GrisClaroUTEC!45,rounded corners=4bp] (196,86) rectangle (742,338);
\Texto{208}{106}{Calibri}{16}{000000}{\bfseries PC2 del robot}
\node[bloqueL,minimum width=124bp,minimum height=58bp] (cam) at (100,170) {Cámara\\RealSense D435i\\[-1bp]{\fontsize{11bp}{12bp}\selectfont USB, en la cabeza}};
\node[bloque,minimum width=200bp,minimum height=74bp,align=left] (srv) at (330,170) {\bfseries \texttt{camara\_servidor.py}\ \ {\mdseries (root)}\\[1bp]{\fontsize{12bp}{14bp}\selectfont IR izquierda 640×480 a 30 fps,}\\[-2bp]{\fontsize{12bp}{14bp}\selectfont emisor apagado; vídeo en :5000}};
\node[bloqueC,minimum width=200bp,minimum height=74bp,align=left] (seg) at (608,170) {\bfseries \texttt{seguidor\_linea.py}\ \ {\mdseries (unitree)}\\[1bp]{\fontsize{12bp}{14bp}\selectfont percepción → estimación → control}\\[-2bp]{\fontsize{12bp}{14bp}\selectfont → supervisor; registro de la tirada}};
\node[bloqueG,minimum width=200bp,minimum height=58bp] (loco) at (608,282) {Control de marcha Unitree\\[-1bp]{\fontsize{11bp}{12bp}\selectfont FSM 201 (alto nivel)}};
\node[bloqueL,minimum width=150bp,minimum height=50bp] (mando) at (850,282) {Mando del operador\\[-1bp]{\fontsize{11bp}{12bp}\selectfont L2 + B: parada}};
\node[bloqueL,minimum width=150bp,minimum height=58bp] (port) at (850,170) {Portátil\\[-1bp]{\fontsize{11bp}{12bp}\selectfont despliegue y análisis}};
\draw[flecha] (cam) -- (srv);
\draw[flecha] (srv) -- node[rotulo,above=1bp] {ZMQ local} (seg);
\draw[flecha] ($(seg.south)+(-50,0)$) -- node[rotulo,left=2bp] {\texttt{Move} a 20 Hz} ($(loco.north)+(-50,0)$);
\draw[flecha] ($(loco.north)+(50,0)$) -- node[rotulo,right=2bp] {IMU, FSM} ($(seg.south)+(50,0)$);
\draw[flecha] (mando) -- (loco);
\draw[flecha,dashed] (port) -- node[rotulo,above=1bp] {WiFi} (seg);
\Nota{208}{306}{270}{DDS de Unitree por \texttt{eth0}: \texttt{rt/lowstate} y LocoClient (RPC)}
% Por qué
\foreach \x/\tit/\txt in {36/{Dos procesos}/{La cámara exige sudo y el DDS de Unitree falla bajo sudo: la cámara va en un proceso root aparte y pasa cada fotograma por ZMQ local (1 ms).},
  337/{El lazo, en el PC2}/{El portátil va por WiFi, con 10–130 ms de ping: solo despliega con rsync, lanza por ssh y analiza los registros.},
  638/{Si el programa calla, se para}/{Cada Move dura 1 s y se manda a 20 Hz. La percepción tarda 3–6 ms por fotograma y el RPC de Move, 0.5 ms.}}{
  \Recuadro{\x}{354}{286}{150}
  \Acento{\x}{354}{504}
  \Parrafo{\x+18}{380}{256}{Calibri}{18}{21}{000000}{\bfseries\tit}
  \Parrafo{\x+18}{406}{256}{Calibri}{15}{18.5}{4E4B4C}{\txt}
}
\end{Lienzo}
\end{frame}

% ── III · Plan de trabajo (sección 8)
% Plan de trabajo según los hitos de la sección 8, con lo que se hizo y cuándo.
\begin{frame}[plain]
\begin{Lienzo}
\Cabecera{III · Plan de trabajo (sección 8)}{Plan por hitos: lo exigido y lo hecho}
\Tabla{36}{80}{p{150bp}p{232bp}p{410bp}p{96bp}}{%
\FCab
\CabI{150}{25}{13.5}{Fase (sección 8)} & \CabI{232}{25}{13.5}{Entregable exigido} & \CabI{410}{25}{13.5}{Lo que hicimos} & \CabI{96}{25}{13.5}{Cuándo}\\\hline
\FGris
\CelNI{150}{45}{13}{Hito 1 · Preparación y medida} & \CelI{232}{45}{13}{Interfaces por escrito; altura, inclinación y zona ciega medidas} & \CelI{410}{45}{13}{Interfaces en \texttt{docs/PLAN.md}; calibración con marcas a 1 y 4 m: 1.66 m, 54.1°, suelo visible de 0.10 a 4.07 m} & \CelI{96}{45}{13}{3 y 5 oct}\\\hline
\FBlanca
\CelNI{150}{45}{13}{Conjunto de datos} & \CelI{232}{45}{13}{Tres recorridos con imágenes e IMU; respuesta al escalón de vyaw} & \CelI{410}{45}{13}{12 recorridos del nivel 1 llevados con \texttt{wasd}; escalones a vx 0 y 0.2. Después, 5 del nivel 2 y 4 del 3} & \CelI{96}{45}{13}{3 y 5 oct; 6–7 nov}\\\hline
\FGris
\CelNI{150}{45}{13}{Hito 2 · Percepción sin robot} & \CelI{232}{45}{13}{Línea superpuesta, \% de detección, tiempo por fotograma y pruebas} & \CelI{410}{45}{13}{99 \% de detección donde se ve la línea (12 recorridos); 3.1–3.4 ms por fotograma en el PC2; 13 pruebas sintéticas} & \CelI{96}{45}{13}{5 oct}\\\hline
\FBlanca
\CelNI{150}{45}{13}{Hito 3 · Control y supervisor en simulacro} & \CelI{232}{45}{13}{Órdenes con el signo correcto; el vigilante para al tapar la cámara} & \CelI{410}{45}{13}{Simulador en lazo cerrado (14/14 con la peor marcha). En el robot: signos bien, PARADA a los 3.0 s con la cámara tapada} & \CelI{96}{45}{13}{5–6 oct}\\\hline
\FGris
\CelNI{150}{45}{13}{Pruebas en el robot por niveles} & \CelI{232}{45}{13}{De menor a mayor velocidad; cada fallo se analiza antes de repetir} & \CelI{410}{45}{13}{36 tiradas registradas; la primera de cada nivel, a media escala; cada cambio, probado antes en el simulador} & \CelI{96}{45}{13}{8 oct – 10 nov}\\\hline
\FBlanca
\CelNI{150}{45}{13}{Métricas y cierre} & \CelI{232}{45}{13}{Métricas de la sección 9; informe, código y registros} & \CelI{410}{45}{13}{\texttt{metricas.py} por tirada; \texttt{RESULTADOS.md} con cada iteración; informe 6.1–6.4 en preparación} & \CelI{96}{45}{13}{nov}\\\hline
}
% Línea de tiempo (octubre y noviembre, con el eje cortado)
\draw[line width=2bp,GrisUTEC] (52,452) -- (420,452);
\draw[line width=2bp,GrisUTEC] (470,452) -- (912,452);
\draw[line width=1.2bp,GrisUTEC] (437,445) -- (445,459) (447,445) -- (455,459);
\Texto{52}{506}{Calibri}{12}{4E4B4C}{\bfseries octubre}
\Texto{470}{506}{Calibri}{12}{4E4B4C}{\bfseries noviembre}
\foreach \x/\fecha/\txt/\arriba in {104/3 oct/{Hito 1, datos y escalones}/1, 190/5 oct/{Hito 2: percepción}/0,
   238/6 oct/{Hito 3: simulacro}/1, 322/8 oct/{1.ª tirada autónoma}/0, 404/10 oct/{Nivel 1 superado}/1,
   482/5 nov/{Giro inicial; datos N2}/0, 580/6 nov/{Nivel 2 superado}/1, 690/7 nov/{1.ª S completa}/0,
   880/10 nov/{Niveles 3 y 4 superados}/1}{
  \fill[CianUTEC,draw=black,line width=0.6bp] (\x,452) circle (5bp);
  \ifnum\arriba=1
    \node[anchor=south,inner sep=0pt,font=\FuenteCalibri\fontsize{12bp}{13bp}\selectfont,align=center] at (\x,440) {\textbf{\fecha}\\\txt};
  \else
    \node[anchor=north,inner sep=0pt,font=\FuenteCalibri\fontsize{12bp}{13bp}\selectfont,align=center] at (\x,464) {\textbf{\fecha}\\\txt};
  \fi
}
\end{Lienzo}
\end{frame}

% ── IV · Desarrollo por hitos
% Hito 1: geometría de la cámara medida (6.1).
\begin{frame}[plain]
\begin{Lienzo}
\Cabecera{IV · Hito 1}{Hito 1: dónde mira la cámara}
% Esquema lateral a 75 bp/m: suelo en y = 292, vertical de la cámara en x = 112
\begin{scope}[shift={(112,292)},x=75bp,y=75bp]
  \fill[CianUTEC,opacity=0.30] (0,1.66) -- (0.10,0) -- (4.07,0) -- cycle;
  \draw[GrisUTEC,line width=1.6bp] (-0.9,0) -- (4.45,0);
  % robot esquemático
  \draw[GrisUTEC,line width=5bp,line cap=round] (-0.14,0.04) -- (-0.10,0.92) (-0.04,0.04) -- (-0.04,0.92);
  \draw[GrisUTEC,line width=9bp,line cap=round] (-0.08,0.98) -- (-0.08,1.50);
  \fill[GrisUTEC] (-0.08,1.62) circle (6bp);
  \fill[black] (0,1.66) circle (3bp);
  \draw[black,line width=0.9bp] (0,1.66) -- (0.10,0);
  \draw[black,line width=0.9bp] (0,1.66) -- (4.07,0);
  \draw[black,dashed,line width=0.9bp] (0,1.66) -- (1.20,0);
  \draw[GrisUTEC,line width=0.6bp] (0,1.66) -- (0.9,1.66);
  \draw[black,line width=0.8bp] (0.55,1.66) arc[start angle=0,end angle=-54.1,radius=0.55];
  \node[rotulo,anchor=west,text=black] at (0.58,1.47) {54.1°};
  \draw[{Stealth[length=5bp]}-{Stealth[length=5bp]},line width=0.8bp] (-0.55,0) -- (-0.55,1.66);
  \node[rotulo,anchor=east,text=black] at (-0.6,0.83) {1.66 m};
  \node[rotulo,anchor=north,text=black] at (0.10,-0.05) {0.10 m};
  \node[rotulo,anchor=north,text=black] at (1.20,-0.05) {1.20 m};
  \node[rotulo,anchor=north,text=black] at (4.07,-0.05) {4.07 m};
  \node[rotulo] at (2.0,0.38) {suelo visible en el IR};
  \node[rotulo,anchor=south west] at (1.2,1.66) {eje óptico a 54.1° bajo la horizontal};
\end{scope}
% Tabla de lo medido
\Tabla{36}{330}{p{168bp}p{120bp}p{136bp}}{%
\FCab
\CabI{168}{24}{13}{Medida} & \CabI{120}{24}{13}{Valor} & \CabI{136}{24}{13}{Método}\\\hline
\FGris
\CelI{168}{27}{12.5}{Altura de la cámara} & \CelN{120}{27}{12.5}{1.66 m} & \CelI{136}{27}{12.5}{cinta (prof.: 1.672)}\\\hline
\FBlanca
\CelI{168}{27}{12.5}{Inclinación} & \CelN{120}{27}{12.5}{54.10°} & \CelI{136}{27}{12.5}{marcas a 1 y 4 m}\\\hline
\FGris
\CelI{168}{27}{12.5}{Suelo visible (IR)} & \CelN{120}{27}{12.5}{0.10–4.07 m} & \CelI{136}{27}{12.5}{desde la vertical}\\\hline
\FBlanca
\CelI{168}{27}{12.5}{Campo de visión IR / color} & \CelN{120}{27}{12.5}{64.7° / 43.4°} & \CelI{136}{27}{12.5}{vertical, intrínsecos}\\\hline
\FGris
\CelI{168}{27}{12.5}{Puntera de los pies} & \CelN{120}{27}{12.5}{−0.02 m} & \CelI{136}{27}{12.5}{de la vertical}\\\hline
}
% Imagen de la calibración
\Doc{484}{84}{440}{330}{calibracion_ir_distancias.png}
\Pie{484}{430}{440}{IR de la calibración: distancias del suelo (amarillo), eje de la cámara y ±30 cm (verde), marca y barra (magenta) y puntos de la línea (azul).}
\Recuadro{484}{462}{440}{44}
\Acento{484}{462}{506}
\Parrafo{500}{479}{418}{Calibri}{14}{16.5}{000000}{\textbf{Se usa el IR}: ve más suelo que el color; el emisor, solo al calibrar.}
\end{Lienzo}
\end{frame}

% Caracterización de la marcha (6.3): respuesta al giro, avance real y balanceo.
\begin{frame}[plain]
\begin{Lienzo}
\Cabecera{IV · Hito 1}{Cómo responde la marcha}
\Figura{36}{84}{430}{235}{fig_escalon}
\Pie{36}{336}{430}{Escalones de vyaw de ±0.3 rad/s andando a vx 0.2 (\texttt{escalon\_20261003\_081312}): orden y giro medido por la IMU.}
\Recuadro{36}{362}{430}{144}
\Acento{36}{362}{506}
\Cuerpo{52}{383}{404}{13.5}{%
\Sub{Lo que cambia en el control}\Esp{4}
\Vin{Pure pursuit desde el centro de giro (14 cm tras la cámara) y sobre la dirección real de avance.}\Esp{3}
\Vin{Compensación fija de la deriva: −0.026 rad/s.}\Esp{3}
\Vin{Punto adelantado a 0.6 m: 0.4 s de retardo son 8 cm a 0.2 m/s.}}
\CifraBarra[30]{516}{86}{408}{0.30–0.45 s}{retardo efectivo del giro andando; ganancia \textasciitilde0.95 y 90 \% en 0.8–1.2 s}
\CifraBarra[30]{516}{190}{408}{×1.14, a 7.5° a la izquierda}{el robot avanza más rápido de lo mandado (0.228 m/s con vx 0.2) y en diagonal respecto del eje de la cámara}
\CifraBarra[30]{516}{294}{408}{1.2–1.7°/s}{de deriva del rumbo a la izquierda andando recto; andando no baja de \textasciitilde0.18 m/s}
\CifraBarra[30]{516}{398}{408}{1.0 → 3.2 cm}{oscilación de la línea al corregir el balanceo con el roll y el pitch de la IMU: empeora. Con el yaw de la IMU, el ángulo pasa de \textasciitilde1° a 0.04–0.13°}
\end{Lienzo}
\end{frame}

% Hito 2: percepción desarrollada sobre los conjuntos de datos, sin mover el robot (6.2).
\begin{frame}[plain]
\begin{Lienzo}
\Cabecera{IV · Hito 2}{Hito 2: la percepción, sin mover el robot}
\Doc{36}{82}{259}{380}{percepcion_ejemplo.jpg}
\Pie{36}{478}{280}{Línea, punto adelantado y barra a 1.89 y 1.00 m; a la derecha, la vista desde arriba.}
% Etapas
\foreach \i/\y/\tit/\txt in {
  1/84/{Vista desde arriba}/{1 cm por píxel, 0.2–3.0 m por delante y ±1 m; geometría calibrada y postura lenta de la IMU.},
  2/158/{Franja de 5 cm}/{Contraste a los dos lados: el canto de un mueble no cuenta. Las dos polaridades.},
  3/232/{Umbral adaptativo}/{6 veces el ruido robusto de cada franja, con un mínimo de 6 niveles.},
  4/306/{Cadena desde el robot}/{La línea empieza bajo el robot; recta o parábola con reajuste local hasta 0.8 m.},
  5/380/{Barra y esquina}/{Cinta transversal donde acaba la línea: a los dos lados, barra; a uno, esquina.}}{
  \Numero{336}{\y}{28}{\i}
  \Parrafo{376}{\y+12}{270}{Calibri}{16}{19}{000000}{\textbf{\tit}}
  \Parrafo{376}{\y+31}{270}{Calibri}{14}{16.5}{4E4B4C}{\txt}
}
\Recuadro{336}{450}{310}{56}
\Acento{336}{450}{506}
\Parrafo{352}{471}{288}{Calibri}{14}{16.5}{000000}{\textbf{Confianza} = filas con línea × contraste × calidad del ajuste.}
% Cifras
\CifraBarra[34]{682}{86}{242}{99 \%}{de detección donde la referencia ve la línea (12 recorridos)}
\CifraBarra[34]{682}{186}{242}{3.1–3.4 ms}{por fotograma en el PC2 (p95 6.1–6.4 ms); a 30 fps hay 33 ms}
\CifraBarra[34]{682}{286}{242}{0.2–2.0 cm}{de diferencia con la referencia; 0.05–0.53° en el ángulo}
\Tabla{670}{376}{p{126bp}p{128bp}}{%
\FCab
\CabI{126}{22}{12.5}{Sombra ×0.25} & \CabI{128}{22}{12.5}{Alcance de la línea}\\\hline
\FGris
\CelI{126}{24}{12.5}{Umbral fijo (15)} & \CelI{128}{24}{12.5}{1.21 m}\\\hline
\FBlanca
\CelNI{126}{24}{12.5}{Adaptativo} & \CelNI{128}{24}{12.5}{2.99 m}\\\hline
}
\Nota{670}{464}{254}{Sombra sintética sobre los datos (6.2.2): el adaptativo sigue la línea.}
\end{Lienzo}
\end{frame}

% Estimación (6.4.1–6.4.4) y ley de control (6.3.3–6.3.6).
\begin{frame}[plain]
\begin{Lienzo}
\Cabecera{IV · Hito 3}{Estimación y ley de control}
\Cuerpo{36}{100}{392}{14.5}{%
\Sub{Estimación: la línea en el suelo, no en la imagen}\Esp{5}
\Vin{Cada medida se guarda en un marco fijo al suelo, girada con el yaw de la IMU de 10 ms antes del fotograma.}\Esp{4}
\Vin{Sin medida, la línea se desplaza con la estima de avance: rumbo por la IMU y 1.14·vx a 7.5°.}\Esp{4}
\Vin{La barra y la esquina valen tras 5 avistamientos y se siguen contando a ciegas.}\Esp{4}
\Vin{Una medida a más de 25 cm o 25° de la línea prevista no se usa.}}
\CifraBarra[30]{50}{346}{378}{1.5–5.7 cm}{de error de la línea prevista tras 2 s a ciegas, sobre los 12 recorridos reales (menos de 0.3°)}
% Esquema de pure pursuit (200 bp/m a escala 0.78: L = 0.6 m). Ocupa de y = 76 (línea, recortada) a
% y = 314 (círculo y rótulo de L), por encima del recuadro de las ecuaciones (y = 328).
\begin{scope}[shift={(672,220)},x=0.78bp,y=-0.78bp]
  \begin{scope}
    \clip (-200,-185) rectangle (200,200);
    \draw[GrisUTEC,line width=3bp] (60,80) .. controls (70,-40) and (40,-150) .. (-40,-220);
  \end{scope}
  \node[rotulo,anchor=east] at (54,72) {línea};
  \draw[GrisMedio,dashed,line width=0.8bp] (0,0) circle (120);
  \draw[GrisMedio,line width=0.6bp] (0,0) -- (-120,0);
  \node[rotulo,anchor=south] at (-66,-2) {L = 0.6 m};
  \fill[GrisUTEC] (0,0) circle (4bp);
  \node[rotulo,anchor=north,text=black] at (0,8) {centro de giro};
  \draw[GrisUTEC,line width=1.2bp] (0,0) -- (0,-34);
  \fill[black] (0,-34) circle (3bp);
  \node[rotulo,anchor=east,text=black] at (-12,-38) {cámara};
  \draw[-{Stealth[length=7bp]},line width=1.5bp,GrisMedio] (0,0) -- (0,-150);
  \draw[-{Stealth[length=7bp]},line width=1.8bp,black] (0,0) -- (-21.9,-166.6);
  \node[rotulo,anchor=east,text=black] at (-24,-140) {avance real (7.5° del eje)};
  \coordinate (P) at (34.2,-115);
  \draw[line width=1.4bp,CianUTEC!80!black] (0,0) -- (P);
  \fill[CianUTEC,draw=black,line width=0.6bp] (P) circle (5bp);
  \node[rotulo,anchor=west,text=black] at (44,-120) {punto adelantado};
  \draw[line width=0.9bp,black] (-5.9,-44.6) arc[start angle=-97.5,end angle=-73.4,radius=45];
  \node[rotulo,text=black] at (9,-66) {$\alpha$};
\end{scope}
\Recuadro{462}{328}{462}{100}
\Acento{462}{328}{428}
\Parrafo{478}{356}{440}{Arial}{15}{20}{000000}{$\omega_{\mathrm{cmd}} = \dfrac{2\,v\,\sin\alpha}{L} + \omega_{\mathrm{sesgo}}$\quad $v = 1.14\,v_x$,\ $\omega_{\mathrm{sesgo}} = -0.026$ rad/s}
\Parrafo{478}{396}{440}{Arial}{14}{18}{000000}{$v_x = 0.4\cdot\text{escala}\cdot\dfrac{1}{1+0.5|\kappa|}\cdot h(\text{confianza})\cdot\cos\alpha$}
\Tabla{462}{436}{p{110bp}p{352bp}}{%
\FCab
\CabI{110}{22}{12.5}{L (simulador)} & \CabI{352}{22}{12.5}{Error máximo de los pies en la S, media escala}\\\hline
\FGris
\CelI{110}{20}{12.5}{1.0 m / 0.8 m} & \CelI{352}{20}{12.5}{15.7 cm / 9.6 cm (10/14 con la peor marcha)}\\\hline
\FCian
\CelNI{110}{20}{12.5}{0.6 m} & \CelNI{352}{20}{12.5}{5.6 cm (14/14 con la peor marcha)}\\\hline
}
\end{Lienzo}
\end{frame}

% Supervisor: máquina de estados explícita, vigilante y condiciones de seguridad (sección 10).
\begin{frame}[plain]
\begin{Lienzo}
\Cabecera{IV · Hito 3}{Supervisor, vigilante y seguridad}
\tikzset{estado/.style={bloque,minimum width=132bp,minimum height=40bp,font=\FuenteCalibri\bfseries\fontsize{14bp}{14bp}\selectfont},
         transicion/.style={rotulo,text=black,fill=white,inner sep=1.5bp}}
\node[estado,fill=GrisClaroUTEC!55] (esp) at (110,108) {ESPERA};
\node[estado,fill=CianUTEC,draw=black] (seg) at (296,226) {SEGUIMIENTO};
\node[estado,fill=GrisClaroUTEC!55] (per) at (482,108) {LÍNEA PERDIDA};
\node[estado,fill=GrisClaroUTEC!55] (esq) at (296,352) {GIRO ESQUINA};
\node[estado,fill=CianUTEC!45] (fin) at (110,446) {FIN};
\node[estado,fill=GrisUTEC,text=white] (par) at (482,446) {PARADA};
\draw[flecha] (esp) -- node[transicion,sloped,above=1bp,pos=0.54] {5 medidas buenas} (seg);
\draw[flecha] ($(seg.north east)+(-14,0)$) -- node[transicion,sloped,above=1bp,pos=0.45] {0.3 s sin medida} ($(per.south west)+(10,0)$);
\draw[flecha] ($(per.south)+(-20,0)$) to[bend left=18] node[transicion,below right=0bp] {recuperada} ($(seg.east)+(0,6)$);
\draw[flecha] (per) -- node[transicion,right=2bp] {4.5 s} (par);
\draw[flecha] ($(seg.south)+(-30,0)$) to[bend right=10] node[transicion,left=2bp] {esquina (nivel 4)} ($(esq.north)+(-30,0)$);
\draw[flecha] ($(esq.north)+(30,0)$) to[bend right=10] node[transicion,right=2bp] {90° y línea nueva} ($(seg.south)+(30,0)$);
\draw[flecha] (seg.west) -| node[transicion,right=2bp,pos=0.72] {barra} (fin.north);
\draw[flecha] (esq.south east) -- node[transicion,sloped,below=1bp] {sin línea} (par.north west);
\draw[flecha] (esp.west) -| (26,494) -- (482,494) -- (par.south);
\node[transicion,anchor=west,align=left] at (30,300) {sin línea\\5 s};
\Nota{196}{94}{214}{Desde cualquier estado: PARADA si salta el vigilante.}
% Seguridad
\Recuadro{602}{82}{322}{422}
\Acento{602}{82}{504}
\Cuerpo{620}{104}{294}{13.5}{%
\Sub{Paradas (StopMove) del vigilante}\Esp{3}
\Vin{Línea perdida más de 4.5 s.}
\Vin{Roll o pitch por encima de 20°.}
\Vin{Motor en fallo o FSM distinta de 201.}
\Vin{Cámara sin fotogramas 0.5 s; \texttt{lowstate} viejo 0.5 s.}
\Vin{Excepción, Ctrl+C o 120 s de tirada.}\Esp{8}
\Sub{Sección 10 en el código}\Esp{3}
\Vin{Topes: vx 0.4 m/s, vy 0.2 m/s, vyaw 0.5 rad/s; primeras tiradas a media escala.}
\Vin{Para andar hay que escribir \texttt{SEGUIR}; la tecla ESPACIO deja el robot quieto de pie.}
\Vin{Cada Move dura 1 s: si el programa calla, el robot para.}
\Vin{Solo alto nivel: sin \texttt{ReleaseMode}, \texttt{rt/lowcmd} ni \texttt{rt/arm\_sdk}.}\Esp{8}
\Sub{En la pista}\Esp{3}
\Vin{Operador con L2\,+\,B listo; franja de 1.5 m libre.}}
\end{Lienzo}
\end{frame}

% Hito 3: simulacro en el robot, simulador en lazo cerrado y pruebas automáticas.
\begin{frame}[plain]
\begin{Lienzo}
\Cabecera{IV · Hito 3}{Hito 3: simulacro, simulador y pruebas}
\Tabla{36}{82}{p{196bp}p{212bp}}{%
\FCab
\CabI{196}{24}{13}{Simulacro (sin ningún Move)} & \CabI{212}{24}{13}{Resultado en el robot}\\\hline
\FGris
\CelI{196}{38}{12.5}{Signos: alineado, girado ±25°, desplazado ±30 cm} & \CelI{212}{38}{12.5}{vyaw con el signo esperado en todos}\\\hline
\FBlanca
\CelI{196}{38}{12.5}{Tapar la cámara 5 s} & \CelI{212}{38}{12.5}{LÍNEA PERDIDA a 0.3 s y PARADA a 3.0 s}\\\hline
\FGris
\CelI{196}{38}{12.5}{Sin línea al empezar} & \CelI{212}{38}{12.5}{PARADA a los 5 s}\\\hline
\FBlanca
\CelI{196}{38}{12.5}{Parar el servidor de cámara} & \CelI{212}{38}{12.5}{PARADA en 0.54 s}\\\hline
\FGris
\CelI{196}{38}{12.5}{Recorrido hasta la barra, con el mando} & \CelI{212}{38}{12.5}{Barra vista de 2.88 a 0.31 m; FIN}\\\hline
}
\Recuadro{36}{328}{408}{176}
\Acento{36}{328}{504}
\Cuerpo{52}{350}{384}{13.5}{%
\Sub{La primera ronda destapó tres fallos}\Esp{3}
\Vin{Reflejos del techo que borraban 50–60 cm de cinta: ahora un hueco se salta si lo de más allá sigue la línea.}
\Vin{Falsos en el borde de la vista que retrasaban la PARADA a \textasciitilde13 s.}
\Vin{En el tramo final a ciegas la velocidad estimada caía: ahora se cuenta con la medida al ver la barra.}}
\Figura{460}{82}{464}{196}{fig_simulacion}
\Pie{460}{296}{464}{Lazo cerrado simulado a escala 1 con el modelo de la marcha medido, pistas del PDF y salidas a ±10°.}
\CifraBarra[30]{476}{336}{200}{14/14}{tiradas simuladas con éxito con todo lo peor de la marcha a la vez}
\CifraBarra[30]{710}{336}{214}{107}{pruebas automáticas: percepción, estimación, control, supervisor y lazo}
\Recuadro{460}{432}{464}{72}
\Acento{460}{432}{504}
\Parrafo{476}{452}{440}{Calibri}{14}{17}{000000}{\textbf{Percepción real sobre imágenes sintéticas} de las pistas: así se probaron la curva, la S, el hueco, la sombra y la esquina antes de montarlas.}
\end{Lienzo}
\end{frame}

% ── V · Resultados por nivel
% Registro de tiradas por fecha: abre los resultados.
\begin{frame}[plain]
\begin{Lienzo}
\Cabecera{V · Resultados}{36 tiradas en el robot, por fecha}
\Tabla{36}{82}{p{80bp}p{58bp}p{84bp}p{72bp}p{66bp}p{528bp}}{%
\FCab
\CabI{80}{24}{13}{Fecha} & \Cab{58}{24}{13}{Nivel} & \Cab{84}{24}{13}{Escala} & \Cab{72}{24}{13}{Tiradas} & \Cab{66}{24}{13}{FIN} & \CabI{528}{24}{13}{Qué se probó}\\\hline
\FGris
\CelI{80}{36}{12.5}{8 oct} & \Cel{58}{36}{12.5}{1} & \Cel{84}{36}{12.5}{0.5} & \Cel{72}{36}{12.5}{2} & \Cel{66}{36}{12.5}{2} & \CelI{528}{36}{12.5}{Primeras tiradas autónomas, salida casi alineada}\\\hline
\FBlanca
\CelI{80}{36}{12.5}{10 oct} & \Cel{58}{36}{12.5}{1} & \Cel{84}{36}{12.5}{0.5 y 1} & \Cel{72}{36}{12.5}{3} & \Cel{66}{36}{12.5}{3} & \CelI{528}{36}{12.5}{Salidas giradas ±10° y escala 1}\\\hline
\FGris
\CelI{80}{36}{12.5}{5 nov} & \Cel{58}{36}{12.5}{1} & \Cel{84}{36}{12.5}{0.5 y 1} & \Cel{72}{36}{12.5}{4} & \Cel{66}{36}{12.5}{4} & \CelI{528}{36}{12.5}{Girar primero en el sitio (umbrales de 12° y 8°)}\\\hline
\FBlanca
\CelI{80}{36}{12.5}{6 nov} & \Cel{58}{36}{12.5}{2} & \Cel{84}{36}{12.5}{0.5 y 1} & \Cel{72}{36}{12.5}{7} & \Cel{66}{36}{12.5}{7} & \CelI{528}{36}{12.5}{Curva de 1 m; velocidad medida tras frenar; reajuste local de la línea hasta 0.8 m}\\\hline
\FGris
\CelI{80}{36}{12.5}{7 nov} & \Cel{58}{36}{12.5}{3} & \Cel{84}{36}{12.5}{0.5} & \Cel{72}{36}{12.5}{3} & \Cel{66}{36}{12.5}{0} & \CelI{528}{36}{12.5}{La S del equipo; pausa para girar}\\\hline
\FBlanca
\CelI{80}{36}{12.5}{7 nov} & \Cel{58}{36}{12.5}{3} & \Cel{84}{36}{12.5}{1} & \Cel{72}{36}{12.5}{7} & \Cel{66}{36}{12.5}{1} & \CelI{528}{36}{12.5}{Escala 1; pausa con la línea perdida; muestreo, rumbo y vx tras la pausa}\\\hline
\FGris
\CelI{80}{36}{12.5}{10 nov} & \Cel{58}{36}{12.5}{3} & \Cel{84}{36}{12.5}{1} & \Cel{72}{36}{12.5}{6} & \Cel{66}{36}{12.5}{6} & \CelI{528}{36}{12.5}{Búsqueda en arco; hueco de 40 cm; sombra en el tramo final y en toda la zona}\\\hline
\FBlanca
\CelI{80}{36}{12.5}{10 nov} & \Cel{58}{36}{12.5}{4} & \Cel{84}{36}{12.5}{0.5 y 1} & \Cel{72}{36}{12.5}{4} & \Cel{66}{36}{12.5}{4} & \CelI{528}{36}{12.5}{Esquina; parada 0.3 m antes y giro a 0.25 rad/s}\\\hline
\FCian
\CelNI{80}{30}{13}{Total} & \CelN{58}{30}{13}{} & \CelN{84}{30}{13}{} & \CelN{72}{30}{13}{36} & \CelN{66}{30}{13}{27} & \CelNI{528}{30}{13}{Además: 47 simulacros, 21 recorridos grabados, 7 calibraciones y 5 escalones}\\\hline
}
\Nota{36}{474}{888}{Datos en \texttt{datos/tirada\_<fecha>\_<hora>\_nivel<N>\_escala<E>/}: fotogramas, medidas, órdenes, \texttt{lowstate}, imágenes IR y \texttt{resumen.json}. Las fechas de las carpetas son las del reloj del robot.}
\end{Lienzo}
\end{frame}

% Nivel 1: recta de 4 m con salida desalineada.
\begin{frame}[plain]
\begin{Lienzo}
\Cabecera{V · Resultados}{Nivel 1: recta de 4 m con salida desalineada}
\Foto{36}{82}{140}{249}{n1_lejos}
\Foto{186}{82}{140}{249}{n1_barra}
\Doc{36}{342}{168}{126}{parada_020241_escala1.jpg}
\Parrafo{214}{356}{112}{Calibri}{13}{15.5}{4E4B4C}{\textbf{Fig.} Los pies sobre la barra al parar (020241, escala 1).}
\BotonVideo{36}{478}{nivel1.mp4}{Vídeo nivel 1}
\Fila{346}{82}{578}{70}{PROBLEMA}{Saliendo girado a la izquierda, el avance en diagonal (7.5°) y el retardo alejan los pies: \textbf{16.7 cm} a los \textasciitilde3 s (015424, salida a −11°).}
\Fila{346}{160}{578}{70}{CAMBIO}{\textbf{Girar primero en el sitio} si al salir |α| > 12°. Con salidas parecidas, los pies se alejan 0.6 cm en vez de 7.0 (095105 frente a 050643).}
\Fila{346}{238}{578}{70}{RESULTADO}{\textbf{9 de 9 tiradas en la barra} sin intervención (8 oct – 5 nov), a media escala y a escala 1. La parada prevista (+10 cm) acierta en ±4 cm.}
\Tabla{346}{322}{p{78bp}p{60bp}p{132bp}p{82bp}p{116bp}p{110bp}}{%
\FCab
\CabI{78}{24}{13}{Tirada} & \Cab{60}{24}{13}{Escala} & \CabI{132}{24}{13}{Salida (línea)} & \Cab{82}{24}{13}{Recorrido} & \Cab{116}{24}{13}{Pies: máximo} & \Cab{110}{24}{13}{Parada prevista}\\\hline
\FGris
\CelI{78}{25}{12.5}{050643} & \Cel{60}{25}{12.5}{0.5} & \CelI{132}{25}{12.5}{−4.6 cm, −1.0°} & \Cel{82}{25}{12.5}{17.2 s} & \Cel{116}{25}{12.5}{11.2 cm} & \Cel{110}{25}{12.5}{+9.1 cm}\\\hline
\FBlanca
\CelI{78}{25}{12.5}{015424} & \Cel{60}{25}{12.5}{0.5} & \CelI{132}{25}{12.5}{−10.1 cm, −11.2°} & \Cel{82}{25}{12.5}{19.1 s} & \Cel{116}{25}{12.5}{16.7 cm} & \Cel{110}{25}{12.5}{+9.8 cm}\\\hline
\FGris
\CelI{78}{25}{12.5}{015841} & \Cel{60}{25}{12.5}{0.5} & \CelI{132}{25}{12.5}{+1.0 cm, +10.0°} & \Cel{82}{25}{12.5}{19.4 s} & \Cel{116}{25}{12.5}{4.3 cm} & \Cel{110}{25}{12.5}{+10.0 cm}\\\hline
\FCian
\CelNI{78}{25}{12.5}{020241} & \CelN{60}{25}{12.5}{1} & \CelNI{132}{25}{12.5}{−7.0 cm, −2.7°} & \CelN{82}{25}{12.5}{11.9 s} & \CelN{116}{25}{12.5}{11.1 cm*} & \CelN{110}{25}{12.5}{+9.0 cm}\\\hline
}
\Nota{346}{484}{578}{*En el arranque; después de los 3 s, 2.3 cm. Pies: centro de los pies respecto de la línea, medido con la cámara. Parada: puntera respecto del borde cercano de la barra.}
\end{Lienzo}
\end{frame}

% Nivel 2: recta, curva de 90° (radio ~1 m) y recta.
\begin{frame}[plain]
\begin{Lienzo}
\Cabecera{V · Resultados}{Nivel 2: curva de 90° de radio 1 m}
\Foto{36}{82}{290}{182}{n2_curva}
\Foto{36}{274}{290}{182}{n2_barra}
\BotonVideo{36}{470}{nivel2.mp4}{Vídeo nivel 2}
\Parrafo{196}{487}{130}{Calibri}{13}{15}{4E4B4C}{Tirada 104456, escala 1}
\Fila{346}{82}{578}{60}{PROBLEMA}{La cinta rayada de peligro junto a la curva tenía más contraste que la línea: se tapó (el PDF pide 1.5 m libres).}
\Fila{346}{150}{578}{60}{PROBLEMA}{Las dos primeras tiradas pararon 17–25 cm antes de la barra: el final a ciegas se contaba con la velocidad de mientras frenaba.}
\Fila{346}{218}{578}{60}{CAMBIO}{Velocidad medida 1.5 s después de frenar y \textbf{reajuste local de la línea hasta 0.8 m}: la parábola de toda la vista sesgaba el ángulo.}
\Figura{346}{290}{330}{170}{fig_nivel2_ajuste}
\Pie{346}{476}{330}{Error del ángulo al acercarse a la curva, frente a la IMU (3 recorridos).}
\Recuadro{692}{290}{232}{214}
\Chip{692}{290}{RESULTADO}
\CifraBarra[30]{716}{326}{200}{7 de 7}{acaban en FIN sin intervención; las 5 posteriores al arreglo, sobre la barra}
\CifraBarra[30]{716}{420}{200}{18.5–20.3 s}{a escala 1, con la parada prevista a +9…+10 cm de la barra}
\end{Lienzo}
\end{frame}

% Nivel 3 (1/2): la S de la pista del equipo y los cambios hasta completarla.
\begin{frame}[plain]
\begin{Lienzo}
\Cabecera{V · Resultados}{Nivel 3: la S con dos semicírculos de 1 m}
\Figura{36}{82}{560}{250}{fig_nivel3_progreso}
\Pie{36}{348}{560}{Cuánto del 2.º semicírculo recorre cada tirada (de 180°, por la IMU), en orden y con el cambio que se probaba. El 1.º semicírculo lo completan todas.}
\Imagen{618}{82}{306}{229.5}{imagenes/fotos/ir_n3_segundo.jpg}
\Pie{618}{330}{306}{Lo que ve la cámara al perderla (104116): la cinta cruza de lado y el detector, un punto por fila, no sigue más allá de \textasciitilde65°.}
% Cambios
\foreach \i/\x/\y/\tit/\txt in {
  1/36/394/{Pausa para girar}/{Si el giro no cabe a la velocidad mínima y |α| > 25°, para y gira en el sitio.},
  2/36/452/{Escala 1}/{El doble de tope de giro (0.5 rad/s) para la curva de 1 m.},
  3/338/394/{Pausa con la línea perdida}/{Si queda a más de 15° (o la cinta a más de 18° de la marcha), para y gira hacia ella.},
  4/338/452/{Muestreo y salida de la pausa}/{La línea medida se prolongaba hacia atrás; tras una pausa, vx ≤ 0.16 durante 2 s.},
  5/640/394/{Búsqueda en arco}/{Si al acabar la pausa sigue sin verla, avanza girando 0.25 rad/s hacia el mismo lado.}}{
  \Numero{\x}{\y}{24}{\i}
  \Parrafo{\x+32}{\y+11}{252}{Calibri}{15}{17}{000000}{\textbf{\tit}}
  \Parrafo{\x+32}{\y+28}{252}{Calibri}{13}{15}{4E4B4C}{\txt}
}
\fill[CianUTEC] (640,452) rectangle (924,506);
\Parrafo{652}{472}{264}{Calibri}{15}{18}{000000}{\textbf{Resultado:} de 1 de 10 tiradas a 6 de 6 en la barra, en 38–41 s.}
\end{Lienzo}
\end{frame}

% Nivel 3 (2/2): interrupción de 40 cm y sombra.
\begin{frame}[plain]
\begin{Lienzo}
\Cabecera{V · Resultados}{Nivel 3: hueco de 40 cm y sombra}
\Foto{36}{82}{290}{163}{n3_s}
\Pie{36}{262}{290}{La S con el hueco en la unión (043803).}
\BotonVideo{36}{284}{nivel3.mp4}{Vídeo nivel 3}
\Imagen{36}{326}{142}{106.5}{imagenes/fotos/ir_luz.jpg}
\Imagen{184}{326}{142}{106.5}{imagenes/fotos/ir_oscuro.jpg}
\Pie{36}{450}{290}{La cámara IR en el mismo punto, con luz (042852) y con la zona a oscuras (044157): ajusta sola la exposición.}
\Tabla{346}{82}{p{70bp}p{196bp}p{78bp}p{82bp}p{80bp}p{72bp}}{%
\FCab
\CabI{70}{24}{13}{Tirada} & \CabI{196}{24}{13}{Condición} & \Cab{78}{24}{13}{Puntera} & \Cab{82}{24}{13}{Recorrido} & \Cab{80}{24}{13}{Sin línea} & \Cab{72}{24}{13}{Pausas}\\\hline
\FGris
\CelI{70}{24}{12.5}{041348} & \CelI{196}{24}{12.5}{S sin hueco} & \Cel{78}{24}{12.5}{+0.01 m} & \Cel{82}{24}{12.5}{38.0 s} & \Cel{80}{24}{12.5}{2.6 s} & \Cel{72}{24}{12.5}{3}\\\hline
\FBlanca
\CelI{70}{24}{12.5}{041539} & \CelI{196}{24}{12.5}{S sin hueco} & \Cel{78}{24}{12.5}{+0.02 m} & \Cel{82}{24}{12.5}{37.9 s} & \Cel{80}{24}{12.5}{3.5 s} & \Cel{72}{24}{12.5}{4}\\\hline
\FGris
\CelI{70}{24}{12.5}{042852} & \CelI{196}{24}{12.5}{Hueco de 40 cm centrado} & \Cel{78}{24}{12.5}{−0.01 m} & \Cel{82}{24}{12.5}{39.6 s} & \Cel{80}{24}{12.5}{2.0 s} & \Cel{72}{24}{12.5}{3}\\\hline
\FBlanca
\CelI{70}{24}{12.5}{043054} & \CelI{196}{24}{12.5}{Hueco de 40 cm centrado} & \Cel{78}{24}{12.5}{+0.03 m} & \Cel{82}{24}{12.5}{40.4 s} & \Cel{80}{24}{12.5}{2.5 s} & \Cel{72}{24}{12.5}{3}\\\hline
\FGris
\CelI{70}{24}{12.5}{043803} & \CelI{196}{24}{12.5}{Hueco + sombra en el tramo final} & \Cel{78}{24}{12.5}{0.00 m} & \Cel{82}{24}{12.5}{40.9 s} & \Cel{80}{24}{12.5}{3.3 s} & \Cel{72}{24}{12.5}{4}\\\hline
\FCian
\CelNI{70}{24}{12.5}{044157} & \CelNI{196}{24}{12.5}{Hueco + toda la zona a oscuras} & \CelN{78}{24}{12.5}{0.00 m} & \CelN{82}{24}{12.5}{41.4 s} & \CelN{80}{24}{12.5}{2.3 s} & \CelN{72}{24}{12.5}{4}\\\hline
}
\Nota{346}{268}{578}{Todas a escala 1 y sin intervención; puntera respecto de la barra al mandar FIN.}
\Figura{346}{290}{400}{215}{fig_nivel3_luz}
\CifraBarra[28]{772}{300}{152}{Sin cambios}{en el código: el hueco y la sombra se pasaron con lo de la S}
\CifraBarra[28]{772}{404}{152}{0.41}{confianza mediana a oscuras en el 2.º semicírculo, frente a 0.55 con luz}
\end{Lienzo}
\end{frame}

% Nivel 4: esquina de 90°, giro sobre el eje y parada adelantada.
\begin{frame}[plain]
\begin{Lienzo}
\Cabecera{V · Resultados}{Nivel 4: esquina de 90° y giro sobre el eje}
\foreach \x/\f/\i/\txt in {36/n4_a/1/{Ve la esquina desde la salida y la confirma},
  261/n4_b/2/{Para con el centro de giro 0.3 m antes},
  486/n4_c/3/{Gira 90° en el sitio con la IMU, a 0.25 rad/s},
  711/n4_d/4/{Sigue la línea nueva hasta la barra}}{
  \Foto{\x}{82}{213}{120}{\f}
  \Numero{\x}{210}{22}{\i}
  \Parrafo{\x+30}{225}{183}{Calibri}{13.5}{15.5}{000000}{\txt}
}
\Fila{36}{262}{430}{70}{PROBLEMA}{Girando 90° en el sitio, el robot se desplaza \textbf{0.30–0.41 m} hacia donde andaba: en la 1.ª tirada, la línea nueva quedó a 0.39 m.}
\Fila{36}{340}{430}{58}{CAMBIO}{Parar 0.3 m antes de la esquina y girar siempre a 0.25 rad/s.}
\Tabla{36}{406}{p{70bp}p{56bp}p{78bp}p{116bp}p{110bp}}{%
\FCab
\CabI{70}{22}{12.5}{Tirada} & \Cab{56}{22}{12.5}{Escala} & \Cab{78}{22}{12.5}{Adelanto} & \Cab{116}{22}{12.5}{Pies tras el giro} & \Cab{110}{22}{12.5}{Desplazamiento}\\\hline
\FGris
\CelI{70}{19.5}{12}{050558} & \Cel{56}{19.5}{12}{0.5} & \Cel{78}{19.5}{12}{0} & \Cel{116}{19.5}{12}{0.39 m pasado} & \Cel{110}{19.5}{12}{0.41 m}\\\hline
\FBlanca
\CelI{70}{19.5}{12}{052855} & \Cel{56}{19.5}{12}{0.5} & \Cel{78}{19.5}{12}{0.3 m} & \Cel{116}{19.5}{12}{0.03 m} & \Cel{110}{19.5}{12}{0.34 m}\\\hline
\FGris
\CelI{70}{19.5}{12}{053202} & \Cel{56}{19.5}{12}{0.5} & \Cel{78}{19.5}{12}{0.3 m} & \Cel{116}{19.5}{12}{0.00–0.02 m} & \Cel{110}{19.5}{12}{0.30–0.33 m}\\\hline
\FCian
\CelNI{70}{19.5}{12}{053907} & \CelN{56}{19.5}{12}{1} & \CelN{78}{19.5}{12}{0.3 m} & \CelN{116}{19.5}{12}{0.00–0.04 m} & \CelN{110}{19.5}{12}{\textasciitilde0.31 m}\\\hline
}
\fill[CianUTEC] (484,262) rectangle (924,298);
\Parrafo{496}{285}{420}{Calibri}{15}{18}{000000}{\textbf{Resultado:} 4 de 4 en la barra; a escala 1, 23.4 s y giro de 89.8°.}
\Figura{484}{306}{440}{166}{fig_nivel4_giro}
\Pie{484}{490}{284}{Pies respecto de la línea nueva (cámara).}
\BotonVideo{774}{478}{nivel4.mp4}{Vídeo nivel 4}
\end{Lienzo}
\end{frame}

% ── VI · Síntesis
% Resumen de resultados con las métricas de la sección 9, por nivel.
\begin{frame}[plain]
\begin{Lienzo}
\Cabecera{VI · Síntesis}{Resumen: métricas de la sección 9}
\Tabla{36}{82}{p{56bp}p{206bp}p{70bp}p{100bp}p{80bp}p{76bp}p{96bp}p{94bp}p{110bp}}{%
\FCab
\CabI{56}{38}{13}{Nivel} & \CabI{206}{38}{13}{Pista montada} & \Cab{70}{38}{13}{Tiradas en FIN} & \Cab{100}{38}{13}{Tirada (escala 1)} & \Cab{80}{38}{13}{Recorrido} & \Cab{76}{38}{13}{Línea vista} & \Cab{96}{38}{13}{Error lateral medio*} & \Cab{94}{38}{13}{Interven\-ciones} & \Cab{110}{38}{13}{Roll / pitch máximos}\\\hline
\FGris
\CelNI{56}{30}{13}{1} & \CelI{206}{30}{13}{Recta de 4 m} & \CelN{70}{30}{13}{9 / 9} & \Cel{100}{30}{13}{020241} & \Cel{80}{30}{13}{11.9 s} & \Cel{76}{30}{13}{92 \%} & \Cel{96}{30}{13}{2.6 cm} & \Cel{94}{30}{13}{0} & \Cel{110}{30}{13}{2.3° / 3.0°}\\\hline
\FBlanca
\CelNI{56}{30}{13}{2} & \CelI{206}{30}{13}{Curva de 90° de radio 1 m} & \CelN{70}{30}{13}{7 / 7} & \Cel{100}{30}{13}{104456} & \Cel{80}{30}{13}{20.3 s} & \Cel{76}{30}{13}{91 \%} & \Cel{96}{30}{13}{4.6 cm} & \Cel{94}{30}{13}{0} & \Cel{110}{30}{13}{2.5° / 1.8°}\\\hline
\FGris
\CelNI{56}{40}{13}{3} & \CelI{206}{40}{13}{S de 1 m, hueco de 40 cm y la zona a oscuras} & \CelN{70}{40}{13}{7 / 16} & \Cel{100}{40}{13}{044157} & \Cel{80}{40}{13}{41.4 s} & \Cel{76}{40}{13}{91 \%} & \Cel{96}{40}{13}{7.2 cm} & \Cel{94}{40}{13}{0} & \Cel{110}{40}{13}{2.5° / 1.7°}\\\hline
\FBlanca
\CelNI{56}{30}{13}{4} & \CelI{206}{30}{13}{Esquina de 90°} & \CelN{70}{30}{13}{4 / 4} & \Cel{100}{30}{13}{053907} & \Cel{80}{30}{13}{23.4 s} & \Cel{76}{30}{13}{84 \%} & \Cel{96}{30}{13}{2.6 cm} & \Cel{94}{30}{13}{0} & \Cel{110}{30}{13}{2.6° / 1.8°}\\\hline
}
\Nota{36}{276}{888}{*Medido por la cámara en su vertical (\textasciitilde la puntera), con \texttt{herramientas/metricas.py}; en las curvas cerradas lo sobrestima, porque la línea cruza de lado. En todas, la parada prevista queda +9…+10 cm pasada la barra; el PDF admite hasta 30 cm.}
\Figura{36}{306}{300}{200}{fig_tiradas}
\CifraBarra[32]{376}{318}{250}{36 tiradas}{registradas en el robot, más 47 simulacros y 21 recorridos grabados}
\CifraBarra[32]{376}{418}{250}{0 intervenciones}{con el mando en las tiradas que llegan a la barra}
\CifraBarra[32]{664}{318}{260}{≤ 3.0°}{de roll y pitch en estas cuatro; el vigilante para en 20°}
\CifraBarra[32]{664}{418}{260}{5–6 ms}{de percepción por fotograma en el robot (p95 10–13 ms)}
\end{Lienzo}
\end{frame}

% Parámetros finales de config/seguidor.yaml, con los que se lograron los resultados.
\begin{frame}[plain]
\begin{Lienzo}
\Cabecera{VI · Síntesis}{Parámetros finales (\texttt{config/seguidor.yaml})}
\Tabla{36}{82}{p{118bp}p{326bp}p{444bp}}{%
\FCab
\CabI{118}{24}{13}{Bloque} & \CabI{326}{24}{13}{Parámetro} & \CabI{444}{24}{13}{Valor y por qué}\\\hline
\FGris
\CelNI{118}{34}{12.5}{Percepción} & \CelI{326}{34}{12.5}{Vista desde arriba; umbral} & \CelI{444}{34}{12.5}{1 cm/px, 0.2–3.0 m y ±1 m; adaptativo, 6 veces el ruido (mínimo 6 niveles)}\\\hline
\FBlanca
\CelNI{118}{34}{12.5}{} & \CelI{326}{34}{12.5}{Reajuste local; hueco máximo} & \CelI{444}{34}{12.5}{0.8 m (sesgo de recta + curva, nivel 2); 0.5 m (la interrupción es de 0.4 m)}\\\hline
\FGris
\CelNI{118}{34}{12.5}{Estimación} & \CelI{326}{34}{12.5}{Avance; velocidad mínima; confirmar} & \CelI{444}{34}{12.5}{1.14·vx a 7.5° a la izquierda; 0.18 m/s; barra y esquina tras 5 avistamientos}\\\hline
\FBlanca
\CelNI{118}{34}{12.5}{Control} & \CelI{326}{34}{12.5}{Punto adelantado; compensación de la deriva} & \CelI{444}{34}{12.5}{L = 0.6 m desde el centro de giro; −0.026 rad/s}\\\hline
\FGris
\CelNI{118}{34}{12.5}{} & \CelI{326}{34}{12.5}{Girar primero; pausa para girar} & \CelI{444}{34}{12.5}{Si |α| > 12° al salir; si el giro no cabe y |α| > 25° (gira hasta la mitad)}\\\hline
\FBlanca
\CelNI{118}{34}{12.5}{} & \CelI{326}{34}{12.5}{Pausa con la línea perdida; búsqueda} & \CelI{444}{34}{12.5}{|α| > 15° o rumbo > 18° (perdida ≥ 0.8 s); arco a vx 0.16 y 0.25 rad/s; vx ≤ 0.16 durante 2 s}\\\hline
\FGris
\CelNI{118}{34}{12.5}{Supervisor} & \CelI{326}{34}{12.5}{Seguir; línea perdida; parada} & \CelI{444}{34}{12.5}{5 medidas con confianza ≥ 0.45; 0.3 s sin medida buena; PARADA a 4.5 s}\\\hline
\FBlanca
\CelNI{118}{34}{12.5}{} & \CelI{326}{34}{12.5}{Fin en la barra} & \CelI{444}{34}{12.5}{Puntera +10 cm pasada la barra, con 0.5 s de frenado (acierta en ±4 cm)}\\\hline
\FGris
\CelNI{118}{34}{12.5}{} & \CelI{326}{34}{12.5}{Esquina} & \CelI{444}{34}{12.5}{Para 0.3 m antes; giro a 0.25 rad/s, tolerancia 3° y 0.6 s estable}\\\hline
\FBlanca
\CelNI{118}{34}{12.5}{Límites} & \CelI{326}{34}{12.5}{Topes y órdenes} & \CelI{444}{34}{12.5}{vx 0.4 m/s, vy 0.2 m/s, vyaw 0.5 rad/s; primeras tiradas a 0.5; Move a 20 Hz}\\\hline
}
\end{Lienzo}
\end{frame}

% Conclusiones y mejoras.
\begin{frame}[plain]
\begin{Lienzo}
\Cabecera{VI · Síntesis}{Conclusiones y mejoras}
\Texto{36}{100}{Calibri}{20}{000000}{\bfseries Conclusiones}
\foreach \i/\y/\tit/\txt in {
  1/114/{Los cuatro niveles, superados}/{Sin intervención y a escala 1 (hasta 0.4 m/s): 11.9 s, 20.3 s, 41.4 s con hueco y a oscuras, y 23.4 s.},
  2/208/{Medir el robot decidió el diseño}/{Avance en diagonal y 14 \% más rápido, 0.3–0.45 s de retardo, no corregir el balanceo paso a paso y el desplazamiento al girar en el sitio.},
  3/302/{Cinco bloques, probados por separado}/{La percepción se hizo sobre datos grabados y el control en el simulador: la primera tirada llegó con el simulacro validado.},
  4/396/{Cada fallo, analizado antes de repetir}/{36 tiradas registradas; cada cambio se probó antes en el simulador y después en el robot.}}{
  \Numero{36}{\y}{30}{\i}
  \Parrafo{78}{\y+15}{420}{Calibri}{17}{20}{000000}{\textbf{\tit}}
  \Parrafo{78}{\y+37}{420}{Calibri}{15}{18}{4E4B4C}{\txt}
}
\Recuadro{530}{82}{394}{422}
\Acento{530}{82}{504}
\Texto{548}{108}{Calibri}{20}{000000}{\bfseries Mejoras}
\Cuerpo{548}{136}{364}{14.5}{%
\Vin{\textbf{Un detector que siga la cinta a cualquier ángulo} (por columnas o por contorno): hoy no pasa de \textasciitilde65° y obliga a pausar en la curva de 1 m.}\Esp{8}
\Vin{\textbf{Explicar el desplazamiento al girar en el sitio} (0.30–0.41 m en 90°) y modelarlo en vez de compensarlo.}\Esp{8}
\Vin{\textbf{Probar vy} para corregir el avance en diagonal (6.3.4).}\Esp{8}
\Vin{\textbf{Ensayar la pista del nivel 3 del PDF} (radios de 1.2 m) y la sesión seguida de los niveles 1 a 4, como pide la evaluación.}\Esp{8}
\Vin{\textbf{Cerrar el informe} con las respuestas 6.1–6.4.}}
\end{Lienzo}
\end{frame}

% Cierre.
\begin{frame}[plain]
\begin{Lienzo}
\Imagen{60}{70}{176}{200}{imagenes/utec_logo2.jpg}
\fill[CianUTEC] (280,80) rectangle (288,262);
\Texto{312}{150}{Calibri}{64}{000000}{\bfseries Gracias}
\Parrafo{314}{196}{600}{Calibri}{19}{24}{4E4B4C}{Miguel Olortegui · Sergio Morales · Axel Cespedes\par Socialización de los retos técnicos – Unitree H1-2 · 13 de noviembre de 2026}
\foreach \x/\f/\n in {36/g1/1, 263/g2/2, 490/g3/3, 717/g4/4}{
  \Foto{\x}{318}{207}{130}{\f}
  \Chip{\x}{318}{NIVEL \n}
}
\Parrafo{36}{486}{888}{Calibri}{14}{17}{4E4B4C}{Código, registros y resultados: \texttt{retos\_capacitacion/seguidor\_linea/} (\texttt{docs/PLAN.md} y \texttt{docs/RESULTADOS.md}).}
\end{Lienzo}
\end{frame}

\end{document}
